% Generated from locale/tr/content/normal-modal-logic/frame-definability/equivalence-S5.tex; do not hand-edit.


\olfileid[tr]{nml}{frd}{es5}

\olsection{Denklik Bağıntıları ve \Log{S5}}

Modal mantık $\Log{S5}$, bütün evrensel çerçevelerde geçerli !!{formula}s
kümesi olarak nitelenir; yani her dünyaya, kendisi de içinde olmak üzere her
dünyadan erişilebilir. Böyle bir durumda $\Box$ zorunluluğa, $\Diamond$ ise
olanağa karşılık gelir: $!A$ \emph{her} dünyada doğruysa $\Box !A$;
\emph{bazı} dünyalarda doğruysa $\Diamond !A$ doğrudur. $\Log{S5}$'in
aynı zamanda bütün yansımalı, simetrik ve geçişli çerçevelerde, yani bütün
\emph{denklik bağıntıları} üzerinde geçerli !!{formula}s olarak da
nitelenebildiği ortaya çıkar.

\begin{defn}
  $W$ üzerinde ikili bir $R$ bağıntısı, ancak ve ancak yansımalı, simetrik ve
  geçişliyse bir \emph{denklik bağıntısıdır}. $W$ üzerindeki bir $R$
  bağıntısı, bütün $u,v\in W$ için $Ruv$ ise \emph{evrenseldir}.
\end{defn}

\Ax{T}, \Ax{B} ve \Ax{4} sırasıyla yansımalı, simetrik ve geçişli
çerçeveleri nitelediğinden, erişilebilirlik bağıntısının bir denklik
bağıntısı olduğu çerçeveler tam olarak bu üç !!{formula}{}ün de geçerli
olduğu çerçevelerdir. Denklik bağıntılarını tanımlarken kullandığımız
koşullar $R$ üzerindeki başka tanıdık koşul birleşimlerine denk olduğundan,
denklik bağıntılarının başka !!{formula} birleşimleriyle de nitelenebildiği
ortaya çıkar.

\begin{prop}\ollabel{prop:equivalences}
  Aşağıdakiler denktir:
  \begin{enumerate}
  \item $R$ bir denklik bağıntısıdır;
  \item $R$ yansımalı ve Öklidyendir;
  \item $R$ seri, simetrik ve Öklidyendir;
  \item $R$ seri, simetrik ve geçişlidir.
  \end{enumerate}
\end{prop}

\begin{proof}
  Alıştırma.
\end{proof}

\begin{prob}
  Aşağıdakileri göstererek \olref[nml][frd][es5]{prop:equivalences}
  önermesini ispatlayın:
  \begin{enumerate}
  \item $R$ simetrik ve geçişliyse Öklidyendir.
  \item $R$ yansımalıysa seridir.
  \item $R$ yansımalı ve Öklidyense simetriktir.
  \item $R$ simetrik ve Öklidyense geçişlidir.
  \item $R$ seri, simetrik ve geçişliyse yansımalıdır.
  \end{enumerate}
  Bunun, koşulların denk olduğunu ispatlamak için neden yeterli olduğunu
  açıklayın.
\end{prob}

\olref{prop:equivalences}, $R$'nin bir denklik bağıntısı olduğu
çerçevelerin modal mantığına (geleneksel olarak $\Log{S5}$ denilen mantığa)
denk bir niteleme verdiği için \olref[prf][prs]{prop:S5} önermesinin
semantik karşılığıdır.

Evrensel bağıntılar ile denklik bağıntıları arasındaki ilişki nedir? Her
evrensel bağıntı bir denklik bağıntısı olsa da her denklik bağıntısının
evrensel olmadığı açıktır. Bununla birlikte bütün evrensel bağıntılarda
geçerli !!{formula}s{}in kümesi, bütün denklik bağıntılarında geçerli
olanların kümesiyle tam olarak aynıdır.

\begin{prop}
  $R$, $W$ üzerinde bir denklik bağıntısı olsun ve her $w\in W$ için $w$'nin
  \emph{denklik sınıfını} $[w]=\{w'\in W:Rww'\}$ olarak tanımlayalım. O
  zaman:
  \begin{enumerate}
  \item $w\in[w]$;
  \item $R$, her $[w]$ denklik sınıfı üzerinde evrenseldir;
  \item Denklik sınıfları topluluğu $W$'yi birbirini dışlayan ve birlikte
    $W$'nin tamamını kapsayan alt kümelere böler.
  \end{enumerate}
\end{prop}

\begin{prop}\ollabel{prop:S5=univ}
  !!^a{formula} $!A$, $R$'nin bir denklik bağıntısı olduğu bütün
  $\mModel{F}=\tuple{W,R}$ çerçevelerinde geçerlidir ancak ve ancak $R$'nin
  evrensel olduğu bütün $\mModel{F}=\tuple{W,R}$ çerçevelerinde geçerlidir.
  Dolayısıyla evrensel çerçevelerin mantığı $\Log{S5}$'tir.
\end{prop}

\begin{proof}
  $W$ üzerindeki evrensel bir $R$ bağıntısının bir denklik bağıntısı olduğu
  hemen doğrulanır. Dolayısıyla $!A$, $R$'nin bir denklik bağıntısı olduğu
  bütün çerçevelerde geçerliyse bütün evrensel çerçevelerde de geçerlidir.
  Öteki yön için karşıt-ters yönde akıl yürütelim: $!B$, $R$'nin $W$
  üzerinde bir denklik bağıntısı olduğu $\tuple{W,R}$ çerçevesine dayanan
  $\mModel{M}=\tuple{W,R,V}$ modelinin bir $w$ dünyasında yanlış olan
  !!a{formula} olsun. Yani $\mSat/{M}{!B}[w]$. Şu şekilde bir
  $\mModel{M}'=\tuple{W',R',V'}$ modeli tanımlayalım:
  \begin{enumerate}
  \item $W'=[w]$;
  \item $R'$, $W'$ üzerinde evrenseldir;
  \item $V'(p)=V(p)\cap W'$.
  \end{enumerate}
  (Dolayısıyla $\mModel{M}'$ modelindeki $W'$ dünyalar kümesi
  \olref{fig:partition} içindeki taralı alanla gösterilir.) $R$ ile $R'$nin
  $W'$ üzerinde aynı olduğu kolayca görülür. Sonra !!{formula}s üzerinde
  tümevarımla, her $!A$ ve her $w'\in W'$ için
  $\mSat{M'}{!A}[w']$ ancak ve ancak $\mSat{M}{!A}[w']$ olduğunu gösterebiliriz
  (bu ifade $W'\subseteq W$ olduğundan anlamlıdır). Özellikle
  $\mSat/{M'}{!B}[w]$ olur ve $!B$, evrensel bir çerçeveye dayanan bir
  modelde yanlışlanır.
\end{proof}

\begin{figure}[t]
  \centering
  \begin{tikzpicture}[node distance=2cm, auto, thick]
    \clip [rounded corners] (0,0) -- (6,0) -- (6,4) -- (0,4) --  cycle;
    \begin{scope}
      \clip (0,0) -- (2,0)  .. controls (1.5,1.5) and (4.5,1.5)
      .. (4,4) -- (0,4) -- (0,0) -- cycle;
      \filldraw[fill=gray!40] (0,4) -- (0,2) .. controls (1.5,1.5)
      and (4.5,1.5) .. (4.5,0) -- (6,0) -- (6,4) -- (0,4) --  cycle;
    \end{scope}
    \draw [name path=line1] (2,0) .. controls (1.5,1.5) and (4.5,1.5) .. (4,4);
    \draw [name path=line2] (0,2)  .. controls (1.5,1.5) and (4.5,1.5) .. (4.5,0);
    \path [name intersections={of=line1 and line2}];
    \draw [rounded corners, use as bounding box, very thick] (0,0)
    -- (6,0) -- (6,4) -- (0,4) --  cycle; 
    \node at (2,3) {$[w]$} ;
    \node at (1,1) {$[u]$} ;
    \node at (3,0.75) {$[v]$} ;
    \node at (4.75,2) {$[z]$} ;
  \end{tikzpicture}
  \caption{$W$'nin denklik sınıflarına bir bölümlemesi.}
  \ollabel{fig:partition}
\end{figure}

